% Chapter 1: Introduction
% Replace this with the actual introduction.

\chapter{Introduction}

\section{Motivation}

Paraguay is a country of 7.4 million people, with a territory spanning 406,752 km². The country is characterized by significant environmental diversity, including the Gran Chaco (a dry forest region in the west) and the Eastern Region (Atlantic Forest remnants, agricultural zones, urban centers).

Despite its importance for the global carbon cycle and biodiversity, Paraguay lacks comprehensive, open-source, AI-driven tools for monitoring land cover change, agricultural yield, indigenous territories, air quality, and wildlife conservation.

This thesis addresses this gap by building \textbf{SatelliteCV-Paraguay}: a unified Python package that ingests multi-temporal satellite imagery and produces six peer-reviewed papers across these critical domains.

\section{Research Questions}

This thesis answers the following questions:

\begin{enumerate}
    \item Can foundation models pre-trained on global satellite data be transferred to Paraguay for land cover and change detection? (P0011 Yvytu, P0100 Yvyra)
    \item Can satellite-based approaches support Paraguay's emerging carbon credit verification needs? (P0100 Yvyra)
    \item Can multi-modal satellite + VLM approaches respect indigenous data sovereignty while identifying territory conflicts? (P0012 Yvy)
    \item Can object detection be applied to wildlife poaching detection in the Gran Chaco? (P0026 Kai)
    \item Can deep learning approaches forecast urban air quality in Asunción? (P0035 Tatakua)
    \item Can soybean yield be predicted using Sentinel-2 + INBIO data? (P0025 Yrupe)
\end{enumerate}

\section{Contributions}

This thesis makes the following contributions:

\begin{itemize}
    \item \textbf{Open-source Python package} (\texttt{satellite-paraguay}) with 6 paper pipelines
    \item \textbf{6 peer-reviewed papers} across remote sensing, agriculture, climate, conservation, and social science
    \item \textbf{Comprehensive data catalog} of 14+ datasets for Paraguay
    \item \textbf{Best practices} for indigenous data sovereignty (CARE Principles)
    \item \textbf{Reproducibility framework} (random seeds, environment capture, DVC)
    \item \textbf{Dashboard} for live monitoring
    \item \textbf{API} for programmatic access
\end{itemize}

\section{Thesis Structure}

This thesis is organized into 11 chapters:

\begin{itemize}
    \item Chapter 1 (this): Introduction
    \item Chapter 2: Literature Review
    \item Chapter 3: Methodology (shared infrastructure)
    \item Chapters 4-9: Six papers
    \item Chapter 10: Integration (the unified platform)
    \item Chapter 11: Conclusions
\end{itemize}
